% GENERATED FILE. Edit canonical lessons, then run npm run generate:books.
% canonical-source-hash: fnv1a64:14a817d284daadee
% canonical-lessons: GU-C47-finger, GU-C47-throat, GU-C47-shoulder, GU-C47-lip, GU-C47-tongue

\chapter{Finger, Throat and Tongue}
\label{ch:gu-finger-throat-and-tongue}
\hlchaptermodality{47}

\begin{chapteropening}
\textbf{By the end of this chapter:} \emph{I can name a finger, the throat, a shoulder, the lips and the tongue.}

Complete the last lesson of chapter 47: I can name a finger, the throat, a shoulder, the lips and the tongue.
\end{chapteropening}

\section[\texorpdfstring{\gu{આંગળી} (ā̃gaḷī)}{આંગળી (ā̃gaḷī)}]{\gu{આંગળી} (ā̃gaḷī) — a finger}

\label{lesson:GU-C47-finger}



\begin{quote}
Before the new one: say the Gujarati for a tooth, then the Gujarati for hair.
\end{quote}

\subsection*{You'll want to know: \gu{આંગળી}}
\textbf{\gu{આંગળી}} (\emph{ā̃gaḷī}) — \textquotedblleft{}a finger\textquotedblright{}. Five on each \textbf{\gu{હાથ}}.

\begin{cousinweb}
\textbf{\gu{આંગળી}} (\emph{āṁgaḷī}) is a tadbhava, a word inherited from Sanskrit and reshaped by sound change. It comes from \emph{aṅguli}, \textquotedblleft{}finger, toe\textquotedblright{} (CDIAL 135).
\end{cousinweb}

\begin{grammarlens}[title={What you've built}]
One more word for finger, throat and tongue.
\end{grammarlens}

\subsection*{Your turn}
\begin{itemize}
\raggedright
  \item \emph{Say it:} \emph{ā̃gaḷī}
  \item \emph{Say it:} \emph{ā̃gaḷī}, once more
  \item \emph{Say it:} say \emph{ā̃gaḷī}
  \item \emph{Recall:} say the Gujarati for a tooth, then the Gujarati for hair, then say \emph{ā̃gaḷī} again
\end{itemize}

\begin{tcolorbox}[breakable,colback=teal!4,colframe=teal!35!black,title={Before you move on}]
What does \gu{આંગળી} mean? (\textquotedblleft{}A finger\textquotedblright{}.) Say it once more.
\end{tcolorbox}
\section[\texorpdfstring{\gu{ગળું} (gaḷũ)}{ગળું (gaḷũ)}]{\gu{ગળું} (gaḷũ) — the throat, the neck}

\label{lesson:GU-C47-throat}



\begin{quote}
Before the new one: say the Gujarati for hair, then the Gujarati for a finger.
\end{quote}

\subsection*{You'll want to know: \gu{ગળું}}
\textbf{\gu{ગળું}} (\emph{gaḷũ}) — \textquotedblleft{}the throat, the neck\textquotedblright{}.

\begin{cousinweb}
\textbf{\gu{ગળું}} (\emph{gaḷũ}) is a tadbhava, a word inherited from Sanskrit and reshaped by sound change. It comes from \emph{gala}, \textquotedblleft{}throat, neck\textquotedblright{} (CDIAL 4070).
\end{cousinweb}

\begin{grammarlens}[title={What you've built}]
One more word for finger, throat and tongue.
\end{grammarlens}

\subsection*{Your turn}
\begin{itemize}
\raggedright
  \item \emph{Say it:} \emph{gaḷũ}
  \item \emph{Say it:} \emph{gaḷũ}, once more
  \item \emph{Say it:} say \emph{gaḷũ}
  \item \emph{Recall:} say the Gujarati for hair, then the Gujarati for a finger, then say \emph{gaḷũ} again
\end{itemize}

\begin{tcolorbox}[breakable,colback=teal!4,colframe=teal!35!black,title={Before you move on}]
What does \gu{ગળું} mean? (\textquotedblleft{}The throat, the neck\textquotedblright{}.) Say it once more.
\end{tcolorbox}
\section[\texorpdfstring{\gu{ખભો} (khabho)}{ખભો (khabho)}]{\gu{ખભો} (khabho) — a shoulder}

\label{lesson:GU-C47-shoulder}



\begin{quote}
Before the new one: say the Gujarati for a finger, then the Gujarati for the throat.
\end{quote}

\subsection*{You'll want to know: \gu{ખભો}}
\textbf{\gu{ખભો}} (\emph{khabho}) — \textquotedblleft{}a shoulder\textquotedblright{}.

\begin{cousinweb}
CDIAL lists \textbf{\gu{ખભો}} (\emph{khabho}) under a reconstructed \emph{*khavaka}, \textquotedblleft{}shoulder-blade\textquotedblright{}, and asks whether it was crossed with \emph{*skambha}, \textquotedblleft{}shoulder-blade, wing\textquotedblright{} (CDIAL 3852, 13640).
\end{cousinweb}

\begin{grammarlens}[title={What you've built}]
One more word for finger, throat and tongue.
\end{grammarlens}

\subsection*{Your turn}
\begin{itemize}
\raggedright
  \item \emph{Say it:} \emph{khabho}
  \item \emph{Say it:} \emph{khabho}, once more
  \item \emph{Say it:} say \emph{khabho}
  \item \emph{Recall:} say the Gujarati for a finger, then the Gujarati for the throat, then say \emph{khabho} again
\end{itemize}

\begin{tcolorbox}[breakable,colback=teal!4,colframe=teal!35!black,title={Before you move on}]
What does \gu{ખભો} mean? (\textquotedblleft{}A shoulder\textquotedblright{}.) Say it once more.
\end{tcolorbox}
\section[\texorpdfstring{\gu{હોઠ} (hoṭh)}{હોઠ (hoṭh)}]{\gu{હોઠ} (hoṭh) — a lip, the lips}

\label{lesson:GU-C47-lip}



\begin{quote}
Before the new one: say the Gujarati for the throat, then the Gujarati for a shoulder.
\end{quote}

\subsection*{You'll want to know: \gu{હોઠ}}
\textbf{\gu{હોઠ}} (\emph{hoṭh}) — \textquotedblleft{}a lip, the lips\textquotedblright{}.

\begin{cousinweb}
\textbf{\gu{હોઠ}} (\emph{hoṭh}) is a tadbhava of Sanskrit \emph{oṣṭha}, \textquotedblleft{}lip\textquotedblright{} (CDIAL 2563). Its \emph{h-} may come from a compound \emph{*adhaoṣṭha}, \textquotedblleft{}lower lip\textquotedblright{}, which K. R. Norman suggested and CDIAL records (CDIAL 245a).
\end{cousinweb}

\begin{grammarlens}[title={What you've built}]
One more word for finger, throat and tongue.
\end{grammarlens}

\subsection*{Your turn}
\begin{itemize}
\raggedright
  \item \emph{Say it:} \emph{hoṭh}
  \item \emph{Say it:} \emph{hoṭh}, once more
  \item \emph{Say it:} say \emph{hoṭh}
  \item \emph{Recall:} say the Gujarati for the throat, then the Gujarati for a shoulder, then say \emph{hoṭh} again
\end{itemize}

\begin{tcolorbox}[breakable,colback=teal!4,colframe=teal!35!black,title={Before you move on}]
What does \gu{હોઠ} mean? (\textquotedblleft{}A lip, the lips\textquotedblright{}.) Say it once more.
\end{tcolorbox}
\section[\texorpdfstring{\gu{જીભ} (jībh)}{જીભ (jībh)}]{\gu{જીભ} (jībh) — the tongue}

\label{lesson:GU-C47-tongue}



\begin{quote}
Before the new one: say the Gujarati for a shoulder, then the Gujarati for a lip.
\end{quote}

\subsection*{You'll want to know: \gu{જીભ}}
\textbf{\gu{જીભ}} (\emph{jībh}) — \textquotedblleft{}the tongue\textquotedblright{}.

\begin{cousinweb}
\textbf{\gu{જીભ}} (\emph{jībh}) is a tadbhava, a word inherited from Sanskrit and reshaped by sound change. It comes from \emph{jihvā}, \textquotedblleft{}tongue\textquotedblright{} (CDIAL 5228).
\end{cousinweb}

\begin{grammarlens}[title={What you've built}]
One more word for finger, throat and tongue.
\end{grammarlens}

\subsection*{Your turn}
\begin{itemize}
\raggedright
  \item \emph{Say it:} \emph{jībh}
  \item \emph{Say it:} \emph{jībh}, once more
  \item \emph{Say it:} say \emph{jībh}
  \item \emph{Recall:} say the Gujarati for a shoulder, then the Gujarati for a lip, then say \emph{jībh} again
\end{itemize}

\begin{tcolorbox}[breakable,colback=teal!4,colframe=teal!35!black,title={Before you move on}]
What does \gu{જીભ} mean? (\textquotedblleft{}The tongue\textquotedblright{}.) Say it once more.
\end{tcolorbox}
